\documentclass[11pt]{article}
\usepackage[utf8]{inputenc}
\usepackage{amsmath,amssymb}
\usepackage{geometry}
\geometry{a4paper, margin=1in}
\usepackage{booktabs}
\usepackage{hyperref}

\title{Documentation of New Calculations for Manuscript Revision}
\author{Antigravity}
\date{\today}

\begin{document}
\maketitle

\section{Overview}
This document records the methodology, parameters, directory organization, and computational rationale for the supplementary calculations performed to address the reviewers' comments for the manuscript ``Coordination-Dependent Hydrogen Adsorption on Co-, Fe-, and Ni-Functionalized Monolayer $\text{CrCl}_3$.

The calculations are organized in the \texttt{crcl3-newcals} directory. Three complementary suites are established:
\begin{enumerate}
    \item \textbf{DFT-D3 (van der Waals) Dispersion Corrections} (\texttt{DFT\_D3/})
    \item \textbf{DFT+U Hubbard Benchmarks} (\texttt{DFT\_U/})
    \item \textbf{Vibrational Frequency Analysis} (\texttt{Freq/})
\end{enumerate}

\section{Methodology and Computational Details}

\subsection{1. Extraction of Adsorption Energy ($\Delta E_{\mathrm{ads}}$) and Clean Substrates}
To rigorously extract the hydrogen adsorption energy:
\begin{equation}
    \Delta E_{\mathrm{ads}} = E(\text{doped-CrCl}_3 + \text{H}) - E(\text{doped-CrCl}_3) - \frac{1}{2} E(\text{H}_2)
\end{equation}
both the hydrogenated slab ($E(\text{doped-CrCl}_3 + \text{H})$) and the corresponding clean substrate ($E(\text{doped-CrCl}_3)$) must be evaluated at the exact same level of theory. Therefore, for both \textbf{DFT-D3} and \textbf{DFT+U}, separate calculations are staged for:
\begin{itemize}
    \item The 6 hydrogenated systems: \texttt{adsorbed/\{Co, Fe, NI\}} and \texttt{embeded/\{Co, Fe, NI\}}.
    \item The 6 clean substrates without H: \texttt{clean/adsorbed/\{Co, Fe, NI\}} and \texttt{clean/embeded/\{Co, Fe, NI\}}.
\end{itemize}
This guarantees consistency, preventing spurious energetic offsets that would arise from mixing plain PBE and corrected total energies.

\subsection{2. DFT-D3 (van der Waals) Setup}
\textbf{Goal:} Address Reviewer 4 (Comment 2) by confirming that long-range dispersion does not alter the relative adsorption trends ($\text{Co} < \text{Fe} < \text{Ni}$) or structural coordination stability.
\begin{itemize}
    \item \textbf{Method:} Grimme D3 zero-damping scheme (\texttt{IVDW = 11}).
    \item \textbf{Ionic Relaxation:} Full relaxation with \texttt{IBRION = 2}, \texttt{NSW = 500}, \texttt{ISIF = 2}, \texttt{EDIFFG = -0.025~eV/\AA}.
    \item \textbf{Spin Polarization:} \texttt{ISPIN = 2} with dopant-specific initial magnetic moments ($3.0\,\mu_{\mathrm{B}}$ for Co, $4.0\,\mu_{\mathrm{B}}$ for Fe, $2.0\,\mu_{\mathrm{B}}$ for Ni, $3.0\,\mu_{\mathrm{B}}$ for Cr).
\end{itemize}

\subsection{3. DFT+U Hubbard Benchmarks}
\textbf{Goal:} Address Reviewer 4 (Comments 1 \& 16) by verifying that the $3d$ electron correlation on Cr does not affect the catalytic ordering governed by TM orbital hybridization.
\begin{itemize}
    \item \textbf{Method:} Rotationally invariant formulation by Dudarev et al. (\texttt{LDAUTYPE = 2}) with $U_{\mathrm{eff}} = 4.0\text{ eV}$ applied exclusively to Cr $3d$ states.
    \item \textbf{INCAR Specifications:}
    \begin{verbatim}
LDAU      = .TRUE.
LDAUTYPE  = 2
LDAUL     = 2 -1 -1 -1   # (2 -1 -1 for clean substrates)
LDAUU     = 4.0 0 0 0    # (4.0 0 0 for clean substrates)
LDAUJ     = 0 0 0 0      # (0 0 0 for clean substrates)
LMAXMIX   = 4            # Mandatory for d-state PAW charge density mixing
    \end{verbatim}
    \item \textbf{Critical Verification:} \texttt{LMAXMIX = 4} is explicitly included to avoid charge mixing truncation at $l=2$.
\end{itemize}

\subsection{4. Vibrational Frequency Calculations}
\textbf{Goal:} Address Reviewer 4 (Comment 4) by obtaining explicit zero-point energy ($\Delta E_{\text{ZPE}}$) and entropic ($-T\Delta S$) corrections for the adsorbed hydrogen.
\begin{itemize}
    \item \textbf{Method:} Finite difference of forces (\texttt{IBRION = 5}, \texttt{NFREE = 2}, displacement \texttt{POTIM = 0.015~\AA}).
    \item \textbf{Precision:} Tight electronic convergence \texttt{EDIFF = 1.0e-07} to eliminate force noise.
    \item \textbf{Constraints:} Selective dynamics in \texttt{POSCAR}. Substrate atoms 1--33 (Cr, Cl, TM) are fixed (\texttt{F F F}), while the adsorbed H atom (atom 34) is allowed to vibrate (\texttt{T T T}). Starting coordinates are initialized directly from fully converged \texttt{CONTCAR} geometries.
\end{itemize}

\section{Summary of Staged Calculations}
A total of \textbf{30 distinct calculations} are staged and ready for cluster execution:
\begin{table}[ht]
\centering
\begin{tabular}{llcc}
\toprule
\textbf{Suite} & \textbf{Configuration} & \textbf{Systems} & \textbf{Total Jobs} \\
\midrule
\texttt{DFT\_D3} & Hydrogenated (adsorbed + embedded) & Co, Fe, Ni & 6 \\
\texttt{DFT\_D3} & Clean Substrates (adsorbed + embedded) & Co, Fe, Ni & 6 \\
\texttt{DFT\_U}  & Hydrogenated (adsorbed + embedded) & Co, Fe, Ni & 6 \\
\texttt{DFT\_U}  & Clean Substrates (adsorbed + embedded) & Co, Fe, Ni & 6 \\
\texttt{Freq}    & Hydrogenated Selective Dynamics & Co, Fe, Ni & 6 \\
\midrule
\textbf{Total}   & & & \textbf{30} \\
\bottomrule
\end{tabular}
\end{table}

\section{Cluster Deployment and Automation}
Execution on the HPC cluster is automated through:
\begin{itemize}
    \item \texttt{DFT\_D3/run\_monitor\_vasp.sh}
    \item \texttt{DFT\_U/run\_monitor\_vasp.sh}
    \item \texttt{Freq/run\_monitor\_vasp.sh}
    \item \texttt{run\_all\_new\_calculations.sh} (root master orchestrator)
\end{itemize}
These scripts check file completeness, track active SLURM jobs to avoid duplication, and assign unique job identifiers (e.g., \texttt{D3\_ads\_Co}, \texttt{U\_cln\_emb\_Fe}, \texttt{Freq\_ads\_Ni}).

\end{document}
